% Part: proof-theory
% Chapter: proof-search
% Section: search-algorithm

\documentclass[../../../include/open-logic-section]{subfiles}

\begin{document}

\olfileid{pt}{ps}{sal}

\olsection{د \Log{G3c} د لټون الګوريتم}

د \Log{G3c} دپاره د لټون يو الګوريتم بيانوو۔ دا يوازېنې
ممکن الګوريتم نۀ دے، او تر ټولو اغېزمن هم نۀ دے۔ خو سم دے:
که سيکوېنټ $\Gamma \Sequent \Delta$ !!a{proof} ولري، نو
بالاخره به ودرېږي او !!a{proof} به ومومي۔ دا بې‌خطا هم دے:
که بې له !!a{proof} موندلو ودرېږي، يا هېڅکله ونۀ درېږي، نو
$\Gamma \Sequent \Delta$ له پيله معتبر نۀ دے۔

د الګوريتم يو ضروري شرط دا دے چې په
$\Gamma \Sequent \Delta$ کښې، يا د ثبوت د لټون په جوړ شوي
سيکوېنټ کښې، هر ناعټومي !!{formula} «راکم» کړي؛ يعنې د
استنتاجي قاعدې په سرچپه کارونه کښې ئې د اصلي فارمول په توګه
وکاروي، او دا کار هر څومره چې اړين وي هغومره تکرار کړي۔ دې
شرط ته \emph{انصاف} وايو۔
\textbf{د سرچينه‌يي شرط سمون OLCMP-191:} د راکمولو وار
ناعټومي فارمولونو ته راځي؛ اټومي فارمول د استنتاجي قاعدې
اصلي مرکب فارمول نۀ دے۔

الګوريتم په پړاوونو کښې کار کوي۔ د هر پړاو پايله د
سيکوېنټونو يوه ونه ده؛ د بل پړاو پايله په هر هغه تر ټولو
پاسني سيکوېنټ کښې د !!a{formula} په راکمولو جوړېږي چې لا
بديهيت نۀ دے۔

په پړاو $0$ کښې له $\Gamma \Sequent \Delta$ څخه
پيل کوو۔ د انصاف د تضمين دپاره په $\Gamma$ او $\Delta$ کښې
!!{formula}s ته شمېرې د شاخصونو په توګه داسې ورکوو چې بېلابېل
!!{formula}s بېلابېلې شمېرې واخلي۔ مثلاً، له کيڼې خوا تر ښي
خوا د !!{formula}s شمېرل په پام کښې ونيسئ۔ شمېرې د پاسنيو
نښو په توګه ليکو۔ که د $!A, !B \Sequent !C$ دپاره !!a{proof}
لټوو، د پړاو~$0$ پايله $!A^1, !B^2 \Sequent !C^3$ ده۔ په
داسې شاخص لرونکي سيکوېنټ کښې د پاسنۍ نښې په توګه راغلے تر
ټولو لوے~$n$ د هغه \emph{اعظمي شاخص} دے (په بېلګه کښې~$3$)۔

که سيکوېنټ کميت ټاکونکي لري، بايد دا هم معلومه وي چې
په ښي اړخ د $\lexists[x][!A(x)]$ يا په کيڼ اړخ د
$\lforall[x][!A(x)]$ د راکمولو دپاره کوم ترمونه وکاروو۔ يوازې
له !!{constant}s $\Obj c_0$، $\Obj c_1$، $\Obj c_2$، \dots
څخه جوړ شوي ترمونه پکار دي، چې په~$\Gamma$ او $\Delta$ کښې
راغلي هر تابع سمبول پکښې کارېدے شي۔ فرض کوو چې د دغو ټولو
ترمونو سېټ په يوه ټاکلې شمېرنه کښې ورکړے شوے، څو مثلاً «هغه
لومړے !!{constant} چې په $\Gamma \Sequent \Delta$ کښې نۀ
راځي» وښيو۔ د الګوريتم د ونې له هر شاخص لرونکي سيکوېنټ سره
د !!{constant}s يو سېټ تړلے وي۔ په پړاو~$0$ کښې سيکوېنټ ته
ورکړے شوے د !!{constant}s سېټ د هغه ټول موجود
!!{constant}s او هم $\{\Obj c_0\}$ رانغاړي؛ نو که سيکوېنټ
هېڅ !!{constant}s ونۀ لري، همدغه وروستے سېټ پاتې کېږي۔
\textbf{د ثابتونو د ثبت سمون OLCMP-194:} پيليز سېټ تل
پاييز ثابت هم رانغاړي، څو په لاندې تکراري کارونه کښې دغه
ثابت د ثبت له سېټه بهر رانۀ شي۔

فرض کړئ چې الګوريتم په پړاو~$n$ کښې د شاخص لرونکو
سيکوېنټونو يوه ونه د !!{constant}s له تړلو سېټونو سره جوړه
کړې ده۔ په پړاو~$n+1$ کښې د هر تر ټولو پاسني سيکوېنټ دپاره
لاندې کار کوو۔

که يو اټومي !!a{formula}~$!A$ د سيکوېنټ په کيڼ او
ښي دواړو اړخونو کښې راځي، يا کيڼ اړخ~$\lfalse$ لري، هېڅ
کار نۀ کوو: داسې سيکوېنټونه په~\Log{G3c} کښې !!{proof}s لري،
او د دغه ځانګړي پاسني سيکوېنټ دپاره نور څۀ پاتې نۀ دي۔
\textbf{د بديهيت د شرط سمون OLCMP-191:} د دواړو خواوو ګډ
فارمول بايد اټومي وي؛ د عمومي مرکب فارمول د هويت ثبوت دلته
په خپله نۀ دے ورګډ شوے۔

که سيکوېنټ هېڅ ناعټومي !!{formula} ونۀ لري، الګوريتم
بند کړئ۔ د $\Gamma \Sequent \Delta$ هېڅ ثبوت نشته۔

کنه، تر ټولو کوچني شاخص~$i$ لرونکے ناعټومي
!!{formula}~$!A$ وټاکئ۔ اړوند پاسنے سيکوېنټ يا
$!A^i, \Pi \Sequent \Lambda$ دے، يا
$\Pi \Sequent \Lambda, !A^i$۔ $k$ دې د سيکوېنټ اعظمي شاخص
او $C$ دې د !!{constant}s هغه سېټ وي چې ورسره تړلے دے۔

$!A^i$ «راکوو»؛ يعنې د~$!A$ د !!{main operator} او
په کيڼ يا ښي اړخ د $!A^i$ د ځاے له مخې د استنتاجي قاعدې
مطابق نوي پاسني سيکوېنټونه جوړوو۔

\begin{enumerate}
  \item $!A \ident !B \land !C$ او په کيڼ اړخ راځي: د
  $(!B \land !C)^i, \Pi \Sequent \Lambda$ پاس يو نوے
  پاسنے سيکوېنټ ورزيات کړئ:
  \[
  \Axiom$!B^{k+1}, !C^{k+2}, \Pi \fCenter \Lambda$
  \RightLabel{\LeftR{\land}}
  \UnaryInf$(!B \land !C)^i, \Pi \fCenter \Lambda$
  \DisplayProof
  \]
  نوي سيکوېنټ ته ورکړے شوے د !!{constant}s سېټ~$C$ دے۔
  \item $!A \ident !B \land !C$ او په ښي اړخ راځي: د
  $\Pi \Sequent \Lambda, (!B \land !C)^i$ پاس دوه نوي
  پاسني سيکوېنټونه ورزيات کړئ:
  \[
  \Axiom$\Pi \fCenter \Lambda, !B^{k+1}$
  \Axiom$\Pi \fCenter \Lambda, !C^{k+1}$
  \RightLabel{\RightR{\land}}
  \BinaryInf$\Pi \fCenter \Lambda, (!B \land !C)^i$
  \DisplayProof
  \]
  نوي سيکوېنټ ته ورکړے شوے د !!{constant}s سېټ~$C$ دے۔
  \textbf{د شاخصونو سمون OLCMP-192، د OLCMP-189 تعقيب:}
  نوې نښې د اوسني اعظمي شاخص څخه پورته ټاکل شوې دي؛ په
  هره يوه مقدمه کښې نوي موردونه بېلابېل شاخصونه اخلي او د
  پاتې موردونو له شاخصونو سره نۀ ټکرېږي۔

  \item د $!A \ident \lnot !B$، $!A \ident !B \lor !C$
  او $!A \ident !B \lif !C$ حالتونه ورته دي او د مشقونو
  دپاره پرېښودل شوي دي۔

  \item $!A \ident \lexists[x][!B(x)]$ او په کيڼ اړخ
  راځي: د $\lexists[x][!B(x)]^i, \Pi \Sequent \Lambda$
  پاس يو نوے پاسنے سيکوېنټ ورزيات کړئ:
  \[
  \Axiom$!B(c)^{k+1}, \Pi \fCenter \Lambda$
  \RightLabel{\LeftR{\lexists}}
  \UnaryInf$\lexists[x][!B(x)]^i, \Pi \fCenter \Lambda$
  \DisplayProof
  \]
  دلته $c$ هغه لومړے !!{constant} دے چې په~$C$ کښې نۀ وي۔
  نوي سيکوېنټ ته ورکړے شوے د !!{constant}s سېټ
  $C \cup \{c\}$ دے۔
  \item $!A \ident \lexists[x][!B(x)]$ او په ښي اړخ راځي:
  د $\Pi \Sequent \Lambda, \lexists[x][!B(x)]$ پاس يو
  نوے پاسنے سيکوېنټ ورزيات کړئ:
  \[
  \Axiom$\Pi \fCenter \Lambda, \lexists[x][!B(x)]^{k+1}, !B(t)^{k+2}$
  \RightLabel{\RightR{\lexists}}
  \UnaryInf$\Pi \fCenter \Lambda, \lexists[x][!B(x)]^i$
  \DisplayProof
  \]
  دلته $t$ هغه لومړے ترم دے چې يوازې د~$C$ !!{constant}s
  لري او تر دې سيکوېنټ لاندې په ښي اړخ د
  $\lexists[x][!B(x)]$ په راکمولو کښې لا نۀ دے کارول شوے
  (يا $\Obj c_0$، که ټول داسې ترمونه کارول شوي وي)۔ نوي
  سيکوېنټ ته ورکړے شوے د !!{constant}s سېټ $C$ دے۔
  \textbf{د رسمي نښو سمون OLCMP-192/193:} نوي موردونه دلته
  هم له پخواني اعظمي شاخص پورته دي؛ د ښي اړخ د موجود کميت
  ټاکنې د قاعدې کاپي شوے کيڼ نوم په سم ښي نوم بدل شوے دے۔
  \item د $!A \ident \lforall[x][!B(x)]$ حالتونه ورته دي
  او د مشقونو دپاره پرېښودل شوي دي۔
\end{enumerate}

که په پړاو $n+1$ کښې هېڅ پاسني سيکوېنټ ته نوے
سيکوېنټ ورزيات نۀ کړو، الګوريتم په برياليتوب درېږي۔ په دې
حالت کښې مو د~$\Gamma \Sequent \Delta$ دپاره !!a{proof}
موندلے، چې د پړاو $n+1$ له ونې څخه د ټولو شاخصونو په لرې
کولو ترلاسه کېږي۔ ټول پاسني سيکوېنټونه يا
$!A, \Pi \Sequent \Lambda, !A$ بڼه لري، چې ګډ فارمول ئې
اټومي دے، يا $\lfalse, \Pi \Sequent \Lambda$ بڼه لري۔ ټول
استنتاجونه د~\Log{G3c} سم استنتاجونه دي۔ دا د بياني
استنتاجونو دپاره څرګنده ده۔ وګورئ چې په هر ګام کښې له
سيکوېنټ سره تړلے د !!{constant}s سېټ~$C$ د هغه ټول موجود
!!{constant}s رانغاړي۔ نو د \LeftR{\lexists} او
\RightR{\lforall} په کارونه کښې د ځانګړي متغير شرط پوره
دے: ځانګړے متغير~$c$ داسې !!a{constant} و چې د نتيجې په
تړلي سېټ~$C$ کښې نۀ و، نو $c$ په نتيجه کښې هم نۀ راتلو۔ لرو:

\begin{prop}
د \Log{G3c} د ثبوت د لټون الګوريتم سم دے: که په برياليتوب
ودرېږي، د پيل سيکوېنټ~$\Gamma \Sequent \Delta$
!!a{proof} شته۔
\end{prop}

\end{document}
